\chapter{Parton Level QED Corrections to Slepton Pair Production}

%%%%%%%%%%%%%%%%%%
%%% Kinematics %%%
%%%%%%%%%%%%%%%%%%
\section{Phase Space Kinematics}
% Introduce d-dimensional phase space for $2\to2$ and $2\to3$ processes, useful kinematic variables like $x, y$ and $z$, and how to transform between quark and slepton \gls{rf}s. Also good place to mention propagator gauge choice and Goldstone bosons?)}
\textcolor{gray}{($d\Pi_{2\to2,3}$ in $d$ dimensions, def. $x$, $y$, $z$, transformation between quark/slepton RF, maybe gauge choice/Goldstone bosons?)}

% Consider a $2 \to n$ scattering process with initial- and final momenta $p_i$ and $k_i$, respectively. Denoting the initial-state energies by $E_{p_i}$, and their velocities by $\vec{p}_{p_i}$, the differential cross section is given by~\cite{ParticleDataGroup:2026aaa}
% \begin{equation}
%     d\hat{\sigma} = \frac{1}{(2E_{p_1})(2E_{p_2})\abs{\vec{v}_{p_1}-\vec{v}_{p_2}}}\abs{\M}^2 d\Pi_{2 \to n},
% \end{equation}
% where $d\Pi_{2 \to n}$ is the $n$-body Lorentz invariant differential phase space, given by
% \begin{equation}
%     d\Pi_{2 \to n} = (2\pi)^d\delta^d \biggl( p_1+p_2 - \sum_{i=1}^n k_i \biggr) \prod_{i=1}^n \frac{d^{d-1}k_i}{(2\pi)^{d-1}}\frac{1}{2E_i}.
% \end{equation}
% The hats above $\hat\sigma$ and future variables are there to remind us that they are \textit{partonic} quantities belonging to the hard scattering of proton constituents. At \gls{nlo} in \gls{qed}, the only hard process we consider is quark--anti-quark scattering, as gluons only appear at higher orders in \gls{qcd}.

Start by considering a $2 \to n$ scattering process with inital- and final momenta $p_i$ and $k_i$, respectively; where particle $k_i$ has mass $m_i$ and energy $E_i$, and the particles $p_i$ are all massless. We further define the \gls{com} energy $\sqrtshat = (p_1+p_2)^2$.
The differential cross section for such a process is~\cite{ParticleDataGroup:2026aaa}
\begin{equation}
    d\hat{\sigma} = \frac{1}{2\sqrtshat} \abs{\M}^2 d\Pi_{2 \to n},
    \label{eq:dsigma}
\end{equation}
where $d\Pi_{2 \to n}$ is the $n$-body Lorentz invariant differential phase space, given in $d$ dimensions by
\begin{equation}
    d\Pi_{2 \to n} = (2\pi)^d\delta^d \biggl( p_1+p_2 - \sum_{i=1}^n k_i \biggr) \prod_{i=1}^n \frac{d^{d-1}k_i}{(2\pi)^{d-1}}\frac{1}{2E_i}.
    \label{eq:dPi_n}
\end{equation}
The hat above $\hat\sigma$ and $\shat$ are there to remind us that they are \textit{partonic} quantities belonging to the hard scattering of proton constituents. It is the hard scattering we will compute perturbatively using Feynman diagrams and \cref{eq:dsigma}. At \gls{nlo} in \gls{qed}, the only hard process we consider is quark--anti-quark scattering, as gluons only appear at higher orders in \gls{qcd}.
\Cref{eq:dPi_n} is somewhat overly general for our purposes, as we will only consider processes with two or three final-state particles. We will explore the two- and three-body phase spaces in more detail in \cref{sec:two_body_phase_space,sec:three_body_phase_space}, but first we turn to a more general discussion of the kinematic variables that will appear in our calculations.

\begin{figure}
    \centering
    \inputtikz{feynman_diagrams/LO}
    \caption{Tree-level diagrams for slepton pair production.}
    \label{fig:M_lo}
\end{figure}

\subsection{Kinematic Variables}
We will consider the $2\to 2(3)$ partonic scattering in which a quark--anti-quark pair annihilates and produces a pair of sleptons, and possibly emits a photon in the process. This process can be summarized as $q\widebar{q} \to \slepton_{iB}^*\slepton_{iA} (+\gamma)$.
The \gls{lo} diagram for this process is shown in \cref{fig:M_lo}, also displaying the definition of momenta for the quarks and sleptons.
In the case of an emitted photon at \gls{nlo}, we denote its momenta by $k = (\omega, \vec{k})$.
Working with scattering at high energies, we assume the quarks to be massless, and use $m_i$ to denote the mass of slepton $k_i$.
It will be useful to define the sum of quark-momenta $p=p_1+p_2$, with $p^2=\shat$, and the sum of slepton-momenta $q=(k_1+k_2)$, with the invariant slepton mass $Q^2\equiv q^2$. Note that $Q^2$ is kinematically constrained by
\begin{equation}
    (m_1+m_2)^2 \leq Q^2 \leq \shat.
    \label{eq:Q2_limits}
\end{equation}

To begin with, we consider the \gls{rf} of the incoming quarks. We are free to choose our coordinate system, and pick the one where the quark moves in the positive $z$-direction, and the photon is emitted at an angle $\theta$ in the $xz$-plane, as shown in \cref{fig:quark_rf}. With this parametrization, we have
\begin{subequations}
\label{eq:momenta_quark_RF}
\begin{align}
    p_1^\mu &= \frac{\sqrtshat}{2} (1, 0, 0, 1),
    \\
    p_2^\mu &= \frac{\sqrtshat}{2} (1, 0, 0, -1),
    \\
    k^\mu &= \frac{\sqrtshat}{2}(1-z) (1, \sin\theta, 0, \cos\theta),
    \label{eq:kmu}
    \\
    q^\mu &= \frac{\sqrtshat}{2}(1-z) 
    \biggl(\frac{1+z}{1-z}, -\sin\theta, 0, -\cos\theta\biggr),
\end{align}
\end{subequations}
where $z=\frac{Q^2}{\shat}$ is the amount of initial-state momenta held by the outgoing sleptons, such that the event of no photon radiation is captured by setting $z=1$. The kinematic constraints on $z$ follows directly from \cref{eq:Q2_limits}:
\begin{equation}
    \frac{(m_1+m_2)^2}{\shat} \leq z \leq 1.
    \label{eq:z_limits}
\end{equation}

We will also need to know the momenta in the sleptons' \gls{rf}. We can transform from the quarks' to the sleptons' \gls{rf} by performing a rotation and a subsequent boost.
A rotation of an angle $\theta$ around the $y$-axis is given by the rotation matrix $\Rot_y(\theta)$; and we define the Lorentz boost such that, if a coordinate $x$ in reference frame $S$ transforms to $x^*$ in $S^*$, where $S^*$ moves with a velocity $\vec{\beta} = \beta\ez$ relative to $S$, then we have $x^* = \Boost_z(\beta)x$. On matrix form, the rotation and boosts are given by
\begin{align}
    \Boost_z(\beta) = \begin{pmatrix}
        \gamma & 0 & 0 & -\gamma\beta \\
        0 & 1 & 0 & 0 \\
        0 & 0 & 1 & 0 \\
        -\gamma\beta & 0 & 0 & \gamma
    \end{pmatrix}, \quad
    \Rot_y(\theta) = \begin{pmatrix}
        1 & 0 & 0 & 0 \\
        0 & \cos\theta & 0 & \sin\theta \\
        0 & 0 & 1 & 0 \\
        0 & -\sin\theta & 0 & \cos\theta
    \end{pmatrix},
\end{align}
where $\gamma = (1-\beta^2)^{-1/2}$ is the Lorentz factor.
The sleptons' \gls{rf} is by definition the frame where $q^*=k_1^*+k_2^*=0$, where we use the asterisk to denote the sleptons' \gls{rf}.
This condition can be reached by first rotating our coordinates in \cref{fig:quark_rf} by an angle $-\theta$, such that $\vec{q}$ points in the negative $z$-direction; and then perform a boost with velocity $-\beta\ez$, where $\beta = \frac{1-z}{1+z}$. The transformation of a coordinate $x^\mu$ in the quarks' \gls{rf} to $x^{*\mu}$ in the sleptons' \gls{rf} can then be summarized as
\begin{align}
    x^{*\mu} &= \Boost_z(-\beta) \Rot_y(-\theta) x^\mu.
    \label{eq:quark_to_slepton_transform}
\end{align}

From \cref{eq:quark_to_slepton_transform} it is fairly straight-forward to find the quark- and photon momenta $p^*$ and $k^*$ in the sleptons' \gls{rf}, and we can parameterize  $k_1^*$ and $k_2^*$ using standard spherical coordinates, as displayed in \cref{fig:slepton_rf}. The four-momenta in the sleptons' \gls{rf} are then
\begin{subequations}
\begin{align}
    p^{*\mu} &= \frac{\shat}{2Q}(1-z) \biggl(\frac{1+z}{1-z}, 0, 0, 1\biggr),
    \\
    k^{*\mu} &= \frac{\shat}{2Q}(1-z) (1, 0, 0, 1),
    \\
    k_i^{*\mu} &= (E_i^*, \pm\kiabs\sin\theta^*\cos\phi^*, \pm\kiabs\sin\theta^*\sin\phi^*, \pm\kiabs\cos\theta^*).
\end{align}
\end{subequations}
Furthermore, using conservation of energy and momentum, we find
\begin{subequations}
\begin{align}
    E_i^* &= \frac{Q^2 \pm (m_1^2-m_2^2)}{2Q},
    \label{eq:Ei}
    \\
    \kiabs &= \frac{\lambda^{\frac{1}{2}}(Q^2, m_1^2, m_2^2)}{2Q},
    \label{eq:kiabs}
\end{align}
\end{subequations}
where $\lambda$ is the K\"all\'en function, defined as
\begin{equation}
    \lambda(a,b,c) = a^2 + b^2 + c^2 - 2ab - 2ac - 2bc,
    \label{eq:kallen}
\end{equation}
for some $a,b,c\in\mathbb{R}$. It is useful to note the bounds of the K\"all\'en function appearing in \cref{eq:kiabs}, which follow from the bounds of $Q^2$ in \cref{eq:Q2_limits}:\footnote{For the upper bound we also set the masses to zero. This imposes a strict inequality, as sleptons are expected to be massive.}
\begin{equation}
    0 \leq \lambda^{\frac{1}{2}}(Q^2, m_1^2, m_2^2) < \shat.
\end{equation}

\begin{figure}
    \centering
    \begin{subfigure}[b]{0.4\linewidth}
        \centering
        \vbox to 4cm {\vfil \inputtikz{figures/quark_rf} \vfil}
        \caption{The quarks' \gls{rf}.}
        \label{fig:quark_rf}    
    \end{subfigure}
    \begin{subfigure}[b]{0.4\linewidth}
        \centering
        \vbox to 4.1cm {\vfil \inputtikz{figures/slepton_rf} \vfil}
        \caption{The sleptons' \gls{rf}.}
        \label{fig:slepton_rf}    
    \end{subfigure}
    \caption{Coordinate system describing the particles' four-momenta in the quarks' and sleptons' \gls{rf}s.}
    \label{fig:rfs}
\end{figure}


\subsection{Two-Body Phase Space}\label{sec:two_body_phase_space}
The two-body phase space is found trivially from \cref{eq:dPi_n}:
\begin{equation}
    d\Pi_{2 \to 2} = (2\pi)^d \delta^d(p-k_1-k_2) \frac{d^{d-1}k_1}{(2\pi)^{d-1}}\frac{1}{2E_1}\frac{d^{d-1}k_2}{(2\pi)^{d-1}}\frac{1}{2E_2},
    \label{eq:dPi_2}
\end{equation}
where $p=p_1+p_2$, and we choose to work in the \gls{com} frame.\footnote{For the $2\to2$ process considered here, the quarks' and sleptons' \gls{rf}, as well as the \gls{com} frame, are identical. Be careful with variable names, however. For instance: $\theta$ and $\theta^*$ are still different, where $\theta^*$ is the angle of $q_1^*$, and $\theta$ is the angle of $k$ (which does not appear in the $2\to2$ case).}
Using the delta function to get rid of the integration over $\vec{k}_2$, as well as the norm of $\vec{k}_1$, the phase space can be reduced to
\begin{equation}
    d\Pi_{2 \to 2} = \frac{\kiabs^{d-3}}{4(2\pi)^{d-2}\sqrtshat} d\Omega^*_{d-1},
\end{equation}
where $d\Omega_d$ is the differential solid angle of the $d$-dimensional unit sphere, and $\kiabs$ was given in \cref{eq:kiabs}.
The differential solid angle can be further simplified using $d\Omega_d = \sin^{d-2}(\theta)d\theta d\Omega_{d-1}$, giving
\begin{equation}
    d\Pi_{2 \to 2} = \frac{\kiabs^{d-3}}{4(2\pi)^{d-2}\sqrtshat} \sin^{d-3}(\theta^*)d\theta^* d\Omega^*_{d-2},
\end{equation}
with $\theta \in [0, \pi)$. For any $2\to 2$ process, working in the \gls{com} frame, there is a rotational symmetry around the collision axis. We can therefore integrate out the remaining solid angle. Using the formula for the surface area of a $d$-dimensional unit sphere~\cite{Schwartz:2014sze},
\begin{equation}
    \Omega_d = \frac{2\pi^{d/2}}{\Gamma\bigl(\frac{d}{2}\bigr)},
\end{equation}
the two-body phase space becomes
\begin{equation}
    d\Pi_{2 \to 2} = \frac{\kiabs^{d-3}}{2(4\pi)^{\frac{d-2}{2}} \Gamma\bigl(\frac{d-2}{2}\bigr) \sqrtshat} \sin^{d-3}(\theta^*)d\theta^*.
\end{equation}


\subsection{Three-Body Phase Space}\label{sec:three_body_phase_space}
Integrating over the three-body phase space is, perhaps unsurprisingly, significantly more cumbersome than the two-body case, and requires some extra care.
From \cref{eq:dPi_n} it follows:
\begin{equation}
    d\Pi_{2\to n} = (2\pi)^d \delta^d(p-k_1-k_2-k) 
    \frac{d^{d-1}k_1}{(2\pi)^{d-1}}\frac{1}{2E_1}
    \frac{d^{d-1}k_2}{(2\pi)^{d-1}}\frac{1}{2E_2}
    \frac{d^{d-1}k}{(2\pi)^{d-1}}\frac{1}{2\omega}.
    \label{eq:dPi_3}
\end{equation}
Working in the \gls{rf} of $p$, we can simplify the differential over $\vec{k}$ by changing to spherical coordinates:
\begin{equation}
    \frac{d^{d-1}k}{2\omega} = \frac{\omega^{d-3}}{2} d\omega d\Omega_{d-1}.
    \label{eq:dPi_3_k}
\end{equation}
To simplify the differentials in $\vec{k}_1$ and $\vec{k}_2$, we will switch to the \gls{rf} of $q=k_1+k_2$. Using the properties of the delta function, along with $q^2=Q^2$, we find
\begin{align}
    \delta^d(p-k_1-k_2-k) 
    \frac{d^{d-1}k_1}{(2\pi)^{d-1}}\frac{1}{2E_1}
    \frac{d^{d-1}k_2}{(2\pi)^{d-1}}\frac{1}{2E_2}
    &= \frac{\kiabs^{d-3}}{4Q} d\Omega^*,
    \label{eq:dPi_3_k12}
\end{align}
with $\kiabs$ defined in \cref{eq:kiabs}.
From \cref{eq:kmu}, after inserting $z=Q^2/\shat$, we have $\omega = \frac{\shat-Q^2}{2\sqrtshat}$, which combined with \cref{eq:dPi_3_k,eq:dPi_3_k12} gives
\begin{equation}
    d\Pi_{2\to3} = \frac{1}{(2\pi)^{2d-3}} \frac{\omega^{d-3}\kiabs^{d-3}}{16Q\sqrtshat} dQ^2 d\Omega_{d-1} d\Omega_{d-1}^*,
\end{equation}
where the integration bounds on $Q^2$ were given in \cref{eq:Q2_limits}.



\section{Factorization of the Cross Section}\label{sec:factorization}

\begin{figure}
    \centering
    \begin{subfigure}[b]{0.45\textwidth}
        \centering
        \inputtikz{feynman_diagrams/factor_left}
        \caption{Initial-state corrections.}
        \label{fig:factor_left}
    \end{subfigure}
    \begin{subfigure}[b]{0.45\textwidth}
        \centering
        \inputtikz{feynman_diagrams/factor_right}
        \caption{Final-state corrections.}
        \label{fig:factor_right}
    \end{subfigure}
    \caption{General Feynman diagrams for the process $q\bar{q} \to \slepton_{iB}^*\slepton_{iA} (+\gamma)$ when considering corrections to either the initial- or final-state part of the \gls{lo} diagram in \cref{fig:M_lo}. The possible addition of an emitted photon with momentum $k$ is represented by light gray external photon lines.}
    \label{fig:M_factored}
\end{figure}

\subsection{Effective Couplings for Initial- and Final-State Diagrams}
The process considered in this thesis is the inclusive production of two sleptons, $q\widebar{q}\to\slepton_{iB}^*\slepton_{iA}(+X)$.
At \gls{lo} there are two diagrams contributing, namely the s-channel diagrams with an intermediate photon or Z-boson, as shown in \cref{fig:M_lo}.
We can simplify the calculation of \gls{nlo} corrections immensely by factoring the \gls{lo} cross section. We will do this in two distinct ways: one where we allow corrections to the initial-state of \cref{fig:M_lo} while keeping the final-state fixed, and one where we allow corrections to the final-state while keeping the initial-state fixed.
In \cref{fig:M_factored} we present the Feynman diagrams for these corrections, 
where the blobs are to be understood as effective couplings, encompassing all \gls{qed} corrections at \gls{nlo} that come from adding photon-lines to the \gls{lo} diagram.
The light gray external photon lines represent the possible emission of an extra photon of momentum $k$, as needed when considering the \textit{inclusive} cross section at \gls{nlo}.
To ease the notation going forward, we will consider such emitted photons as being part of the blobs in \cref{fig:M_factored}, allowing the blobs to represent corrections for both the $2\to2$ and $2\to3$ scattering processes.

As mentioned, the blobs represent \gls{nlo} corrections that add photon lines to the \gls{lo} diagram in \cref{fig:M_lo}. 
Therefore, even at \gls{nlo}, our diagrams will still contain the $q\bar{q}(\gamma/Z)$ vertex, and must either contain the $\slepton\slepton^*(\gamma/Z)$ coupling, or the four-point $\slepton\slepton^*\gamma(\gamma/Z)$ vertex.
We can therefore factor out part of the "blob" coupling, which will be common to all considered diagrams. Applying the Feynman rules in \cref{tab:slepton_rules} gives the following "effective" Feynman rules:
\begin{align}
    \inputtikz[baseline=(v.base)]{rules/factor_left}
    &= \begin{cases}
        -ieQ_q \Gamma^\mu \,,\quad (\gamma)
        \\
        \frac{ie}{\sw\cw} \Gamma^\mu  \Gamma^5 \,,\quad (Z)
    \end{cases} \label{eq:factor_left}
    \\
    \inputtikz[baseline=(v.base)]{rules/factor_right}
    &= \begin{cases}
        ie \delta^{AB} K^\mu \,,\quad (\gamma)
        \\
        \frac{ie}{\sw\cw} \zliab K^\mu \,,\quad (Z)
    \end{cases}, \label{eq:factor_right}
\end{align}
where we have defined the shorthand $\Gamma^5 = \zql(1-\gamma^5)+\zqr(1+\gamma^5)$. $\Gamma^\mu$ and $K^\mu$ represent effective couplings containing the remaining part of the diagrams, such that at \gls{lo} we have $\Gamma^\mu=\gamma^\mu$ and $K^\mu=(k_2-k_1)^\mu$.
Note in particular that $\Gamma^\mu$ has Dirac structure, as it will contain gamma matrices,\footnote{It will, however, not contain $\gamma^5$, as we only consider corrections in \gls{qed}, which is chirally symmetric.} while $K^\mu$ does not. In addition, from the available momenta in the diagrams we can see that $\Gamma^\mu$ depends on $k_1$ and $k_2$ only through $q$, and similarly $K^\mu$ depends on $p_1$ and $p_2$ only through $p$.
One should also note that the chiral factor $\Gamma^5$ in the $q\bar{q}Z$ coupling of \cref{eq:factor_left} has been moved to the right. With each additional photon propagating from one of the quark lines, we get an extra $q\bar{q}\gamma$ coupling, as well as a quark propagator. We therefore always add an even number of gamma matrices, allowing $\Gamma^5$ to safely move past them.
Since the blobs represent corrections both with and without photon emission, $\Gamma^\mu$ and $K^\mu$ may also contain photon polarization vectors.

Before moving further, let us briefly discuss which assumptions we have made, and how they limit the generalization of our procedure.
The perhaps most obvious limitation is that our factorization only works when considering higher-order corrections to either the initial- or final-state of the process. For initial--final-state interferences, which we will consider in \cref{sec:interference_diagrams}, we have to turn back to consider their full diagrams without factorization.
The other assumption we made, which is important for the way we wrote down the effective couplings in \cref{eq:factor_left,eq:factor_right}, is that the \gls{nlo} corrections add photon lines to the \gls{lo} diagram. This does not hold for supersymmetric corrections, such as the one in \cref{fig:M_susy}. This limitation is good to keep in mind, altough we will not consider supersymmetric corrections in this thesis.

\begin{figure}
    \centering
    \inputtikz{feynman_diagrams/susy}
    \caption{Example of a supersymmetric \gls{qed} correction involving a neutralino $\tilde\chi^0$ and two squarks $\tilde{q}$. The factorization described in \cref{sec:factorization} can not be used here as the diagram lacks the $q\bar{q}(\gamma/Z)$ coupling.}
    \label{fig:M_susy}
\end{figure}


\subsection{The Squared Amplitude}
Let us first consider the case where we leave the final-state as it was at \gls{lo}, while allowing corrections to the initial-state. This process is encapsulated in \cref{fig:factor_left}, and using our newfound Feynman rules for the effective coupling in \cref{eq:factor_left}, and remembering that $q^2=Q^2$, we can write down the amplitude. 
To declutter notation we introduce the shorthands $u_1 \equiv u(p_1)$ and $\widebar{v}_2 \equiv \widebar{v}(p_2)$ for the Dirac spinors, and, anticipating the use of dimensional regularization at \gls{nlo}, we work in $d=4-2\varepsilon$ dimensions:
\begin{align}
    \M^i_\gamma &= -\frac{e^2Q_q\mu^{4-d}\delta^{AB}}{Q^2} \widebar{v}_2\Gamma^\mu u_1 (k_2-k_1)_\mu
    \\
    \M^i_Z &= \frac{e^2\mu^{4-d}\zliab}{\sw^2\cw^2}\frac{1}{Q^2-m_Z^2+im_Z\Gamma_Z}\widebar{v}_2\Gamma^\mu\Gamma^5 u_1 (k_2-k_1)_\mu
    \\
    \M^i = \M^i_\gamma+\M^i_Z
    &= -e^2\mu^{4-d}(k_2-k_1)_\mu 
    \\ &\quad \times \biggl\{
        \frac{Q_q\delta^{AB}}{Q^2}\widebar{v}_2\Gamma^\mu u_1 - \frac{\zliab}{\sw^2\cw^2}\frac{1}{Q^2-m_Z^2+im_Z\Gamma_Z}\widebar{v}_2\Gamma^\mu\Gamma^5 u_1
    \biggr\}.
\end{align}
The upper index $i$ is there to remind us that this considers initial-state corrections.
We then compute the spin- and color-averaged squared amplitude using the completeness relations for spinors \rednote{ref?}
\begin{align}
    \overline{|\M^i|^2}
    &= \frac{e^4\mu^{2(4-d)}}{4N} (k_2-k_1)_\mu (k_2-k_1)_\nu \Biggl\{
        \frac{Q_q^2\delta^{AB}}{Q^4}\Tr\Bigl[\ptwoslash\Gamma^\mu \poneslash\widebar\Gamma^\nu \Bigr]
        \nonumber \\
        &\quad - 2\Re\biggl\{ \frac{Q_q\delta^{AB}\zliab}{\sw^2\cw^2}\frac{1}{Q^2\bigl[ Q^2-m_Z^2+im_Z\Gamma_Z \bigr]} \Tr\Bigl[ \ptwoslash\Gamma^\mu\poneslash\widebar\Gamma^\nu\Gamma^5 \Bigr] \biggr\}
        \nonumber \\
        &\quad + \frac{(\zliab)^2}{\sw^4\cw^4}\frac{1}{(Q^2-m_Z^2)^2+m_Z^2\Gamma_Z^2}\Tr\Bigl[ \ptwoslash\Gamma^\mu\poneslash\widebar\Gamma^\nu \Gamma^5\Gamma^5 \Bigr]
    \Biggr\}, \label{eq:Mi_squared_pre_trace}
\end{align}
where we understand $\widebar\Gamma^\mu$ to be $\Gamma^\mu$ with all gamma matrices in reverse order and c-numbers complex conjugated, so as to fulfill $(\widebar{v}_2\Gamma^\mu u_1)^\dagger = \widebar{u}_1 \widebar\Gamma^\mu v_2$. Note also that, in the $2\to3$ case where an extra photon is emitted, there is an implicit sum over photon polarization vectors.

The trace of $\gamma^5$ with an even number of $\gamma^\mu$ reduces to a sum of terms proportional to the levi-civita tensor together with metric tensors, for different permutations of the involved indices:
\begin{equation}
    \Tr[\gamma^{\mu_1}\cdots\gamma^{\mu_n}\gamma^5] = 
    \sum_{\mathclap{ \substack{\text{permutations} \\ \text{of }\mu_1,\ldots,\mu_n}}} 
    \Bigl(\propto\varepsilon^{\mu_1\mu_2\mu_3\mu_4} g^{\mu_5\mu_6}\cdots g^{\mu_{n-1}\mu_n}\Bigr),
\end{equation}
Since the squared amplitude is a scalar quantity, those Lorentz indices must be contracted with Lorentz indices from the remaining expression in \cref{eq:Mi_squared_pre_trace}, whose available tensor quantities (after integrating out $k_1$ and $k_2$, of which the trace is independent) are the four-momenta $p_1^\mu$, $p_2^\mu$, and $q^\mu$;\footnote{$k^\mu$ can be eliminated by momentum conservation: $k=p_1+p_2-q$.} and the metric tensor $g^{\mu\nu}$.
Consider the term with $\varepsilon^{\mu_1\mu_2\mu_3\mu_4}$. Contracting this Levi-Civita tensor with anything that is symmetric in two of those indices yield zero, so the only possible non-zero case is if the tensor it is contracted with contains all the available tensors: $p_{1\mu_1}p_{2\mu_2}q_{\mu_3}g_{\mu_4\mu_5}$. However, since $n$ is even, and there are only metric tensors left to contract with our expression, $g_{\mu_4\mu_5}$ must at some point get contracted with $p_1^{\mu_5}$, $p_2^{\mu_5}$ or $q^{\mu_5}$, resulting in a tensor that is symmetric in at least two of the indices $\mu_1,\ldots,\mu_4$. Therefore, all traces with a single $\gamma^5$ becomes zero when contracted with the remaining part of the cross section.\rednote{Kort ned argumentet!}
Using this, and the fact that $\gamma^5$ squares to unity, we can show that the three terms in \cref{eq:Mi_squared_pre_trace} are actually proportional, and we can rewrite the squared amplitude as
\begin{align}
    \overline{|\M^i|^2} = \frac{16\pi^2\alpha^2\mu^{2(4-d)}}{Q^4}\fqliab(Q^2) \frac{1}{4N} \Tr\Bigl[ \ptwoslash\Gamma^\mu\poneslash\widebar\Gamma^\nu \Bigr] (k_2-k_1)_\mu (k_2-k_1)_\nu, \label{eq:Mi_squared}
\end{align}
where we have defined the effective coupling \rednote{unødvendig å definere på nytt hvis jeg bruker LO-delen}
\begin{align}
    \fqliab(Q^2) &= Q_q^2\delta^{AB}
    - \frac{2Q_q\delta^{AB}\zliab(\zql+\zqr)}{\sw^2\cw^2}\frac{Q^2(Q^2-m_Z^2)}{(Q^2-m_Z^2)^2+m_Z^2\Gamma_Z^2}
    \nonumber\\ &\quad + \frac{2\bigl(\zliab\bigr)^2(\zql^2+\zqr^2)}{\sw^4\cw^4}\frac{Q^4}{(Q^2-m_Z^2)^2+m_Z^2\Gamma_Z^2}.
\end{align}
The result for final-state corrections is derived in the same way, starting from \cref{fig:factor_right}. The squared amplitude is then similar to \cref{eq:Mi_squared}, but with the following differences:
\begin{itemize}
    \item The intermediate vector boson has momentum $p$, instead of $q$, with $p^2=\shat$.
    \item The initial-state coupling reduces to the \gls{lo} case, so $\Gamma^\mu$ becomes $\gamma^\mu$.
    \item We use the effective coupling on the final-state as given in \cref{eq:factor_right}, so we must substitute $(k_2-k_1)^\mu \to K^\mu$.
\end{itemize}
The resulting spin- and color-averaged squared amplitude for final-state corrections is then
\begin{equation}
    \overline{|\M^f|^2} = \frac{16\pi^2\alpha^2\mu^{2(4-d)}}{\shat^2}\fqliab(\shat) \frac{1}{4N} \Tr\Bigl[ \ptwoslash\gamma^\mu\poneslash\gamma^\nu \Bigr] K_\mu K_\nu. \label{eq:Mf_squared}
\end{equation}

\subsection{The Differential Cross Section}
The differential cross section is given by \rednote{ref/cite}
\begin{equation}
    d\hat{\sigma} = \frac{1}{2\shat} \overline{|\M|^2} d\Pi.
\end{equation}
For final-state corrections, as seen in \cref{eq:Mf_squared}, all dependence on phase space momenta are contained in $K_\mu K_\nu$, and we can write the final-state differential cross section as
\begin{equation}
    \frac{d\hat{\sigma}^f}{dQ^2} = \frac{8\pi^2\alpha^2}{\shat^3} \fqliab(\shat) 
    X_f^{\mu\nu}
    L_{\mu\nu},
    \label{eq:dsigmaf_LX}
\end{equation}
where we have defined the tensors
\begin{align*}
    L_{\mu\nu} &= \mu^{4-d} \int d\Pi_L K_\mu K_\nu,
    \\
    X_f^{\mu\nu} &= \frac{\mu^{4-d}}{4N} \Tr\Bigl[ \ptwoslash\gamma^\mu\poneslash\gamma^\nu \Bigr],
\end{align*}
with $d\Pi_L$ defined by extracting the differential $dQ^2$ from \rednote{ref $d\Pi$}, such that
\begin{equation}
    d\Pi = dQ^2 d\Pi_L.
\end{equation}
We can compute $X_f^{\mu\nu}$ using the trace relations for gamma matrices \rednote{ref}, and since it does not acquire any higher-order corrections, we can safely set $d=4$, giving
\begin{align}
    \frac{d\hat{\sigma}^f}{dQ^2} &= \frac{8\pi^2\alpha^2}{N\shat^3}\fqliab(\shat) \biggl(p_{1\mu} p_{2\nu} + p_{1\nu} p_{2\mu} - \frac{\shat}{2}g_{\mu\nu} \biggr) L^{\mu\nu}.
    \label{eq:dsigmaf_tensor}
\end{align}

The squared amplitude for initial-state corrections in \cref{eq:Mi_squared} has a Dirac trace that only depends on $k_1$ and $k_2$ implicitly through $q$, thus we can simplify the differential cross section further in this case by also factoring the phase space.
We consider the phase space for a $2\to3$ scattering process \rednote{ref $d\Pi_{2\to3}$}:\footnote{In the case of no photon emission, $k$ can be trivially integrated out, and the phase space reduces to \rednote{ref. $d\Pi_{2\to2}$}.}
\begin{equation}
    d\Pi = (2\pi)^d\delta^d(k_1+k_2+k-p_1-p_2) \frac{d^{d-1}k}{(2\pi)^{d-1}}\frac{1}{2\omega}\frac{d^{d-1}k_1}{(2\pi)^{d-1}}\frac{1}{2E_1}\frac{d^{d-1}k_2}{(2\pi)^{d-1}}\frac{1}{2E_2}.
\end{equation}
We multiply by the factors $\frac{d^d q}{(2\pi)^d} (2\pi)^d \delta^d(k+q-p_1-p_2)$ and $\delta(q^2-Q^2)dQ^2$, which both integrate to unity. The phase space then becomes
\begin{align}
    d\Pi &= (2\pi)^d \delta^d(q+k-p_1-p_2) \frac{d^dq}{(2\pi)^d}\frac{d^{d-1}k}{(2\pi)^{d-1}}\frac{1}{2\omega} \delta(q^2-Q^2)dQ^2
    \nonumber \\ &\quad \times (2\pi)^d \delta^d(k_1+k_2-q) \frac{d^{d-1}k_1}{(2\pi)^{d-1}}\frac{1}{2E_1}\frac{d^{d-1}k_2}{(2\pi)^{d-1}}\frac{1}{2E_2}.
\end{align}
We can use $\delta(q^2-Q^2)$ to get rid of the integral over the zeroth component of $d^dq$, leaving us with the factored phase space
\begin{equation}
    d\Pi = \frac{1}{2\pi} d\Pi_X d\Pi_H dQ^2,
    \label{eq:phase_space_factored}
\end{equation}
where we have defined the hadronic- and sleptonic phase spaces:
\begin{align}
    d\Pi_H &= (2\pi)^d \delta^d(q+k-p_1-p_2) \frac{d^{d-1}q}{(2\pi)^{d-1}}\frac{1}{2q^0}\frac{d^{d-1}k}{(2\pi)^{d-1}}\frac{1}{2\omega},
    \label{eq:phase_space_H}
    \\
    d\Pi_X &= (2\pi)^d \delta^d(k_1+k_2-q) \frac{d^{d-1}k_1}{(2\pi)^{d-1}}\frac{1}{2E_1}\frac{d^{d-1}k_2}{(2\pi)^{d-1}}\frac{1}{2E_2}.
    \label{eq:phase_space_L}
\end{align}
Using this, we can write the differential cross section for initial-state corrections as
\begin{equation}
    \frac{d\hat{\sigma}^i}{dQ^2} = \frac{4\pi\alpha^2}{\shat Q^4}\fqliab(Q^2)
    H^{\mu\nu} X^i_{\mu\nu},
    \label{eq:dsigmai_HX}
\end{equation}
where we have defined the following tensors:
\begin{align}
    H^{\mu\nu} &= \frac{\mu^{4-d}}{4N}\int d\Pi_H \Tr\Bigl[ \ptwoslash\Gamma^\mu\poneslash\widebar\Gamma^\nu \Bigr],
    \\
    X_i^{\mu\nu} &= \mu^{4-d} \int d\Pi_X (k_2-k_1)^\mu(k_2-k_1)^\nu.
    \label{eq:Ximunu}
\end{align}
As with $X_f^{\mu\nu}$ in \cref{eq:dsigmaf_LX}, $X_i^{\mu\nu}$ does not receive any higher-order corrections, so we can solve it in $d=4$ dimensions.
With the available tensors, it is clear that it must have the form\footnote{The factor $Q^2$ is inserted so that $X_1(Q)$ and $X_2(Q)$ are both dimensionless quantities.}
\begin{equation}
    X_i^{\mu\nu} = X_1(Q) Q^2 g^{\mu\nu} + X_2(Q) q^\mu q^\nu.
    \label{eq:X_decomp}
\end{equation}
Contracting with $g_{\mu\nu}$ and $q_\mu q_\nu$ we get the following coupled equations
\begin{align}
    g_{\mu\nu}X_i^{\mu\nu} &= Q^2 \bigl[ 4 X_1(Q) + X_2(Q) \bigr],
    \\
    q_\mu q_\nu X_i^{\mu\nu} &= Q^4 \bigl[ X_1(Q) + X_2(Q) \bigr],
\end{align}
which can be solved to give the following expression for $X_1(Q)$:
\begin{align}
    X_1(Q) &= -\frac{1}{3Q^4} \Bigl[ q_\mu q_\nu X_i^{\mu\nu} - Q^2 g_{\mu\nu} X_i^{\mu\nu} \Bigr].
    \label{eq:X1_decomp}
\end{align}
We could in principle also calculate $X_2(Q)$, but it is not necessary as the second term in \cref{eq:X_decomp} disappears when contracting with $H^{\mu\nu}$ by the Ward identity. \rednote{Se nærmere på dette. Ikke åpenbart for meg at Ward skal holde når $\gamma^*$ er off-shell}
Inserting \cref{eq:Ximunu} into \cref{eq:X1_decomp}, and using that $q=k_1+k_2$, we can show that (with $d=4$):
\begin{align}
    X_1(Q) &= -\frac{4\kiabs^2}{3Q^2} \int d\Pi_X
    \\
    &= -\frac{\kiabs^3}{3\pi Q^3},
    \label{eq:X1}
\end{align}
where the integral over $d\Pi_X$ was solved using \rednote{ref $d\Pi_{2\to2}$ in $4$ dimensions}, and $\kiabs$ was defined in \rednote{ref}. We can now insert \cref{eq:X_decomp} into \cref{eq:dsigmai_HX} with our expression in \cref{eq:X1} for $X_1(Q)$, resulting in
\begin{equation}
    \frac{d\hat{\sigma}^i}{dQ^2} = -\frac{4\alpha^2 \kiabs^3}{3\shat Q^5}\fqliab(Q^2)
    g_{\mu\nu} H^{\mu\nu}.
    \label{eq:dsigmai_tensor}
\end{equation}

The remaining tensors $H^{\mu\nu}$ and $L^{\mu\nu}$ in \cref{eq:dsigmaf_tensor,eq:dsigmai_tensor} can be rewritten in terms of Feynman diagrams involving an "effective" photon $\gamma^*$. Its Feynman rules are the same as that of the photon, but with some numerical factors left out:\footnote{The same numerical factors that appear in front of $\Gamma^\mu$ and $K^\mu$ in \cref{eq:factor_left,fig:factor_right}.}
\begin{equation}
    \inputtikz[baseline=(v)]{rules/peff_q_qbar} = \gamma^\mu,
    \quad
    \inputtikz[baseline=(v)]{rules/peff_l_l} = (k_2-k_1)^\mu,
    \quad
    \inputtikz[baseline=(v)]{rules/peff_p_l_l} = 2eg^{\mu\nu}.
\end{equation}
From these Feynman rules we have:
\begin{align}
    \varepsilon_\mu \varepsilon_\nu^* H^{\mu\nu} &\equiv \mu^{4-d} \int d\Pi_H \overline{|\M(q\bar{q} \to \gamma^*(+\gamma)|^2},
    \label{eq:H_munu}
    \\
    \varepsilon_\mu \varepsilon_\nu^* L^{\mu\nu} &\equiv \mu^{4-d} \int d\Pi_L \overline{|\M(\gamma^* \to \slepton_{iB}^*\slepton_{iA} (+\gamma)|^2}.
    \label{eq:L_munu}
\end{align}
Using that $-g_{\mu\nu}$ is equal to the sum of polarizations $\varepsilon_\mu\varepsilon^*_\nu$, we can simplify the initial-state cross section in \cref{eq:dsigmai_tensor} even further, giving
\begin{align}
    \frac{d\hat{\sigma}^i}{dQ^2} &= \frac{4\alpha^2 \kiabs^3}{3\shat Q^5}\fqliab(Q^2) \varepsilon_\mu\varepsilon^*_\nu H^{\mu\nu}
    \\
    &= \frac{4\alpha^2 \kiabs^3}{3\shat Q^5}\fqliab(Q^2) 
    \mu^{4-d} \int d\Pi_H \overline{\abs{\M(q\bar{q}\to\gamma^*(+\gamma)}^2}.
    \label{eq:dsigmai_scalar}
\end{align}

To summarize this section: we have laid the groundwork for doing calculations at \gls{nlo}, by simplifying in two major ways:
\begin{itemize}
    \item We have reduced what was originally a $2\to2(3)$ scattering process down to a $2\to1(2)$ scattering for initial-state corrections, given in \cref{eq:dsigmai_scalar}; and $1\to2(3)$ decay for final-state corrections, given in \cref{eq:dsigmaf_tensor}; greatly simplifying the phase space integration.
    \item To begin with, we had to consider two versions of all diagrams, one with an intermediate photon, and one with a Z-boson. This has now been reduced to a single effective photon $\gamma^*$ which incorporates both vector bosons.
\end{itemize}


\section{The Leading Order Cross Section}
We will calculate the cross section at \gls{lo} in two ways, using the initial-state and final-state expressions in \cref{eq:dsigmaf_tensor,eq:dsigmai_tensor}.
Because there is no photon emission at \gls{lo}, these will actually turn out to be the same,\footnote{At least in four spacetime dimensions. Since we set $d=4$ in part of \cref{eq:dsigmaf_tensor,eq:dsigmai_scalar},
the result will be slightly different in the two cases when working in $d=4-2\varepsilon$ dimensions.} but keeping track of $\shat$ and $Q^2$ in the two expressions will be helpful when writing \gls{nlo} calculations in terms of the \gls{lo} result.

\begin{figure}
    \centering
    \begin{subfigure}[t]{0.45\textwidth}
        \centering
        \inputtikz{feynman_diagrams/lo_left}
        \caption{Initial-state case.}
        \label{fig:M_lo_left}
    \end{subfigure}
    \begin{subfigure}[t]{0.45\textwidth}
        \centering
        \inputtikz{feynman_diagrams/lo_right}
        \caption{Final-state case.}
        \label{fig:M_lo_right}
    \end{subfigure}
    \caption{Factored Feynman diagrams for \gls{lo} "correction" to the initial- and final-state differential cross sections in \cref{eq:dsigmai_scalar,eq:dsigmaf_tensor}.}
    \label{fig:M_lo_factored}
\end{figure}

\subsection{The Initial-State Way}
The initial-state \gls{lo} Feynman diagram is shown in \cref{fig:M_lo_left}, giving the amplitude
\begin{align}
    \M^i_0 &= \widebar{v}_2\gamma^\mu u_1 \varepsilon^*_\mu(q).
\end{align}
The spin- and color averaged squared amplitude is then
\begin{align}
    \overline{\abs{\M^i_0}^2} &= \frac{-g_{\mu\nu}}{4N}\Tr\Bigl[ \ptwoslash \gamma^\mu \poneslash \gamma^\nu \Bigr]
    \\
    &= \frac{d-2}{2N}\shat,
\end{align}
where we have used \rednote{ref $\gamma^\mu\gamma^\nu\gamma_\mu$}.
Inserting this into \cref{eq:dsigmai_scalar}, we find
\begin{align}
    \frac{d\hat{\sigma}_i^\mathrm{LO}}{dQ^2} &= \frac{4\alpha^2 \kiabs^3}{3\shat Q^5}\fqliab(Q^2) 
    \mu^{4-d} \frac{d-2}{2N}\shat \int \frac{d^{d-1}q}{(2\pi)^{d-1}}\frac{1}{2q^0} (2\pi)^d \delta^d(q-p_1-p_2)
    \\
    &= \biggl\{ \frac{d-2}{2} \mu^{4-d} \biggr\}
    \frac{8\pi\alpha^2 \kiabs^3}{3N\shat Q^5}\fqliab(Q^2) \delta(1-z)
    \\
    &\equiv \hat{\sigma}_i^{0} \delta(1-z),
    \label{eq:sigma0i}
\end{align}
where the factor in curly brackets reduces to unity in $d=4$ spacetime dimensions.

\subsection{The Final-State Way}
The final-state \gls{lo} Feynman diagram is shown in \cref{fig:M_lo_right}, and we can write down its amplitude:
\begin{align}
    \M^f_0 &= (k_2-k_1)^\mu\varepsilon_\mu(q).
\end{align}
Inserting into \cref{eq:L_munu} we have
\begin{equation}
    L^{\mu\nu} = \mu^{4-d} \int d\Pi_L (k_2-k_1)^\mu (k_2-k_1)^\nu.
    \label{eq:L_munu_lo}
\end{equation}
Using the momenta given in \rednote{ref} we can find the cross section from \cref{eq:dsigmaf_tensor}:
\begin{align}
    \frac{d\hat{\sigma}_f^\mathrm{LO}}{dQ^2} &= \frac{8\pi^2\alpha^2}{N\shat^3}\fqliab(\shat) \biggl(p_{1\mu} p_{2\nu} + p_{1\nu} p_{2\mu} - \frac{\shat}{2}g_{\mu\nu} \biggr) L^{\mu\nu}
    \\
    \hat{\sigma}_f^\mathrm{LO} &= \frac{8\pi^2\alpha^2}{N\shat^3}\fqliab(\shat) \mu^{4-d}
    \int d\Pi \biggl[
        2 \Bigl(p_1 \cdot (k_2-k_1)\Bigr)
        \Bigl(p_2 \cdot (k_2-k_1)\Bigr)
        - \frac{\shat}{2}(k_2-k_1)^2
    \biggr]
    \\
    &= \frac{8\pi^2\alpha^2}{N\shat^3}\fqliab(\shat) \mu^{4-d}
    2\shat \kiabs^2 \int d\Pi \sin^2\theta^*.
    \label{eq:sigmaf_lo_dpi}
\end{align}
The last integral is found using \rednote{ref $d\Pi_{2\to2}$}:
\begin{align}
    \int d\Pi \sin^2\theta^* &= \frac{\pi^{\frac{d-2}{2}}\kiabs^{d-3}}{2\sqrtshat(2\pi)^{d-2}\Gamma(\frac{d-2}{2})}
    \int d\theta^* \sin^{d-1}{\theta^*}
    \\
    &= \frac{\kiabs}{6\pi\sqrtshat}
    \frac{3(d-2)\Gamma\bigl(\frac{d-2}{2}\bigr)}{2(d-1)\Gamma(d-2)} \biggl(\frac{\kiabs^2}{\pi}\biggr)^{\frac{d-4}{2}},
\end{align}
which inserted back into \cref{eq:sigmaf_lo_dpi} gives
\begin{align}
    \hat{\sigma}_f^\mathrm{LO} &= \frac{8\pi^2\alpha^2}{N\shat^3}\fqliab(\shat)
    2\shat \kiabs^2
    \frac{\kiabs}{6\pi\sqrtshat}
    \frac{3(d-2)\Gamma\bigl(\frac{d-2}{2}\bigr)}{2(d-1)\Gamma(d-2)}
    \biggl(\frac{\pi \mu^2}{\kiabs^2}\biggr)^{\frac{4-d}{2}}
    \\
    \frac{d\hat{\sigma}_f^\mathrm{LO}}{dQ^2}
    &= \Biggl\{
        \frac{3(d-2)\Gamma\bigl(\frac{d-2}{2}\bigr)}{2(d-1)\Gamma(d-2)}
        \biggl(\frac{\pi \mu^2}{\kiabs^2}\biggr)^{\frac{4-d}{2}}
    \Biggr\}
    \frac{8\pi\alpha^2 \kiabs^3}{3N\shat^{7/2}}\fqliab(\shat)
    \delta(1-z)
    \\
    &\equiv \hat{\sigma}_f^0 \delta(1-z),
    \label{eq:sigma0f}
\end{align}
where the factor inside the curly brackets reduces to unity in $d=4$ spacetime dimensions.
Note also that if we set $Q^2 = \shat$ and $d=4$, \cref{eq:sigma0i,eq:sigma0f} give the same result, which is expected, as it shouldn't matter \textit{how} we compute the \gls{lo} cross section.


%%%%%%%%%%%%%%
%%%%%%%%%%%%%%
%%%%%%%%%%%%%%
%%%%%%%%%%%%%%




%%%%%%%%%%%
%%% NLO %%%
%%%%%%%%%%%
\section{Next-To-Leading Order Corrections}
\subsection{Initial-State Corrections}
\subsubsection{Virtual Contribution}
\subsubsection{Real Photon Emission}

\subsection{Final-State Corrections}
\subsubsection{Virtual Contribution}
\subsubsection{Real Photon Emission}

\subsection{Initial--Final-State Interference}\label{sec:interference_diagrams}
\subsubsection{Box Diagrams}
\subsubsection{Real Photon Emission}

\section{Summary}
